% ============================================================================
%  SECTION 3 — LE PROBLEME : MINIMISATION CONVEXE COMPOSITE   (Membre 1)
% ============================================================================
\section{Le probleme : minimisation convexe composite}
\label{sec:probleme}

\subsection{Un probleme inverse sous-determine}

Placons-nous dans une situation typique du traitement du signal. On observe un
vecteur de mesures $y \in \R^{n}$, relie a un signal inconnu $x \in \R^{p}$ par
une matrice connue $A \in \R^{n \times p}$ :
\[
  y = A x + \varepsilon,
\]
ou $\varepsilon$ represente un bruit. Le point determinant est que l'on dispose
de \emph{moins de mesures que d'inconnues} ($p > n$) : le systeme est
\emph{sous-determine} et admet une infinite de solutions compatibles avec $y$.
Sans information supplementaire, le probleme est mal pose.

L'information que l'on ajoute est une hypothese de \textbf{parcimonie}
(\emph{sparsity}) : le vrai signal $x$ ne possede que peu de composantes non
nulles. Cette hypothese est realiste --- de nombreux signaux naturels admettent
une representation creuse dans une base adaptee --- et suffit a rendre le
probleme bien pose.

\subsection{Formulation : le LASSO}

Pour retrouver $x$, on minimise un objectif en deux termes, l'un mesurant la
fidelite aux donnees, l'autre favorisant la parcimonie :
\begin{equation}
\label{eq:lasso}
  \boxed{\;
  \min_{x \in \R^{p}} \; F(x)
  = \underbrace{\tfrac{1}{2}\,\norm{A x - y}_2^{2}}_{f(x)\ \text{(lisse)}}
  \;+\; \lambda \underbrace{\norm{x}_1}_{g(x)\ \text{(non lisse)}}
  \;}
\end{equation}
ou $\norm{x}_1 = \sum_{i=1}^{p} |x_i|$ et $\lambda > 0$ regle le compromis entre
les deux termes. Ce probleme est connu sous le nom de \textbf{LASSO}
(\emph{Least Absolute Shrinkage and Selection Operator}). On notera dans toute
la suite
\[
  f(x) = \tfrac{1}{2}\norm{A x - y}_2^{2},
  \qquad
  g(x) = \lambda \norm{x}_1,
  \qquad
  F = f + g .
\]
Le terme $f$ est differentiable, de gradient
\begin{equation}
\label{eq:gradf}
  \nabla f(x) = A\T (A x - y),
\end{equation}
qui est \emph{lipschitzien} de constante $L = \norm{A\T A}_2$ (la plus grande
valeur propre de $A\T A$). Cette constante jouera un role central dans le choix
du pas des algorithmes (sections~\ref{sec:ista} et~\ref{sec:fista}).

\subsection{Le probleme est convexe}

Verifions que~\eqref{eq:lasso} entre bien dans le cadre du cours.

\begin{proposition}
\label{prop:convexe}
La fonction objectif $F = f + g$ est convexe sur $\R^{p}$, et l'ensemble
admissible $\R^{p}$ est convexe. Le probleme~\eqref{eq:lasso} est donc un
probleme d'optimisation convexe.
\end{proposition}

\begin{proof}
On procede terme a terme.
\begin{enumerate}
  \item \emph{$f$ est convexe.} La fonction $f(x) = \tfrac12\norm{Ax-y}_2^2$ est
        quadratique, de matrice hessienne $\nabla^2 f(x) = A\T A$. Pour tout
        $u \in \R^{p}$, $u\T A\T A\, u = \norm{A u}_2^2 \ge 0$, donc $A\T A$ est
        semi-definie positive : $f$ est convexe.
  \item \emph{$g$ est convexe.} La norme $\norm{\cdot}_1$ est une norme, donc
        convexe (inegalite triangulaire et homogeneite) ; multipliee par
        $\lambda > 0$, elle reste convexe
        (proposition~\ref{prop:stabilite}).
  \item \emph{$F$ est convexe.} Comme somme de deux fonctions convexes, $F$ est
        convexe (proposition~\ref{prop:stabilite}).
  \item \emph{L'ensemble admissible.} Il s'agit de $\R^{p}$ tout entier, qui est
        trivialement convexe.
\end{enumerate}
La propriete-or (theoreme~\ref{th:or}) s'applique donc : tout minimum local
de~\eqref{eq:lasso} est global.
\end{proof}

\begin{remarque}[Pas de stricte convexite en general]
\label{rem:pas-stricte}
Conformement a la remarque~\ref{rem:unicite}, $F$ n'est pas strictement convexe
des que $p > n$. En effet, $f$ n'est strictement convexe que si $A\T A$ est
\emph{definie} positive, c'est-a-dire si $A$ est injective ; or, avec plus de
colonnes que de lignes ($p > n$), $A$ possede un noyau non trivial. Quant a
$g = \lambda\norm{\cdot}_1$, elle est convexe mais pas strictement. On ne peut
donc pas garantir l'unicite du minimiseur $\xopt$ : les algorithmes des sections
suivantes convergeront vers une \emph{meme valeur optimale} $\Fopt$, sans que le
point atteint soit necessairement unique.
\end{remarque}

\subsection{Pourquoi les methodes du cours echouent}
\label{subsec:echec}

Le terme $g(x) = \lambda\norm{x}_1$ fait intervenir des valeurs absolues
$|x_i|$. Or la fonction $t \mapsto |t|$ presente un \emph{point anguleux} en
$0$ : sa derivee vaut $-1$ a gauche et $+1$ a droite, et \textbf{n'existe pas en
$0$}. Cette simple observation a des consequences directes :
\begin{itemize}
  \item la \textbf{descente de gradient} requiert $\nabla F$, qui n'est pas
        defini la ou des composantes de $x$ s'annulent --- c'est-a-dire
        precisement la ou se trouve la solution recherchee, puisqu'elle est
        creuse ;
  \item la \textbf{methode de Newton} exige en outre la hessienne de $F$, encore
        moins disponible ;
  \item le \textbf{simplexe} ne traite que les objectifs lineaires, alors que le
        notre est quadratique.
\end{itemize}
Il nous faut donc un outil capable de gerer ce point anguleux sans recourir a la
differentiabilite. Cet outil est l'\emph{operateur proximal}, introduit a la
section suivante.
